\chapter{A deep dive into \mytool}\label{chap:approach}

% --------------------- Section 1 ---------------------

\section{Introduction}\label{sec:approach-introduction}
In Chapter~\ref{chap:running_example}, we presented a practical example showcasing the usage of \mytool{}.
This chapter dives deeper into its design principles and implementation details, elucidating the underlying mechanisms that enable stable specialization through metaprogramming.
The actual implementation is in Rust, but as stated earlier, the concepts can be generalized to other programming languages with similar capabilities.
First, we begin by outlining the core design goals that guided the development of the tool.

\paragraph{Design Goals}
The primary design goals for \mytool{} were centered around achieving effective specialization in stable code while adhering to the principles of safety, usability and performance.
These goals shaped the tool's architecture, macro design and code generation techniques.
The key objectives included:
\begin{itemize}
	\item \textbf{Stable Compatibility}: Ensure that the tool operates entirely within the confines of stable Rust, avoiding any reliance on nightly features or unstable APIs.
	\item \textbf{Type Safety}: Maintain Rust's strong type safety guarantees, ensuring that specialized implementations are correctly typed and validated at compile time.
	\item \textbf{Ease of Use}: Provide a straightforward and intuitive interface for developers to define specialized implementations without excessive boilerplate or complexity.
	\item \textbf{Performance Efficiency}: Generate specialized code that is optimized for performance, minimizing overhead and ensuring that the benefits of specialization are realized in practice.
\end{itemize}

\paragraph{Overview}
As already seen, \mytool{} leverages the macro system to enable specialization.
The three primary components needed by the approach are:
\begin{itemize}
	\item \textbf{\nameref{sec:conditions_definition}}: A macro that annotates specialized implementations, specifying the conditions under which they should be used and generates the necessary boilerplate.
	\item \textbf{\nameref{sec:method_wrapping}}: A second macro that wraps method calls, directing them to the appropriate specialized implementation.
	\item \textbf{\nameref{sec:build_script}}: A way to analyze the code at compile time and collect the necessary information to facilitate specialization.
\end{itemize}

% --------------------- Section 2 ---------------------

\section{Conditions Definition}\label{sec:conditions_definition}
The first step to achieve specialization through metaprogramming is to define various specialized implementations that should be applied when certain conditions are met.
In \mytool{}, this is accomplished using the \inlinerust{#[when(predicate)]} attribute macro, which takes as argument the predicate, which is the condition that must be satisfied for the specialized implementation to be selected.
Please note that if we wanted a simpler drop-in replacement for the unstable specialization feature, we could avoid passing the predicate altogether and just use normal rust syntax for implementations, but we still would need to mark the implementation with some kind of macro to generate the necessary boilerplate.
However, we believe that our approach of using a predicate provides more flexibility and control over the specialization process, allowing for more fine-grained optimizations and better integration with existing codebases.

\paragraph{Predicates}
In our implementation, the macro supports these predicates:
\begin{itemize}
	\item \textbf{Type Equality}: Check if a generic type parameter is of a specific type, e.g.\@ \mbox{\texttt{T = i32}}.
	\item \textbf{Trait Implementation}: Check if a generic type parameter implements a specific trait or multiple traits, e.g.\@ \mbox{\texttt{T:\@ Display}} or \mbox{\texttt{T:\@ Debug + Clone}}.
	\item \textbf{Logical And}: Check if all the inner predicates are satisfied, e.g.\@ \inlinerust{all(predicate1, predicate2, ...)}.
	\item \textbf{Logical Or}: Check if at least one of the provided predicates is satisfied, e.g.\@ \inlinerust{any(predicate1, predicate2, ...)}.
	\item \textbf{Logical Not}: Check if the provided predicate is not satisfied, e.g.\@ \inlinerust{not(predicate)}.
\end{itemize}
These predicates can be combined to form complex conditions, allowing for flexible and powerful specialization rules.

\paragraph{Wildcards}
In addition to the specific predicates, \mytool{} supports a wildcard condition that matches any type.
For instance, the condition \mbox{\texttt{T = Vec<\_>}} matches any vector type, regardless of its element type.
Another useful wildcard condition is \mbox{\texttt{T = \&\_}}, which matches any reference type.
Wildcards are not allowed in trait implementation predicates, as from our point of view they would not provide any useful functionality, for example knowing that a type implements any trait without specifying which one doesn't help much in specializing code.

\paragraph{Nested Conditions}
The macro also supports nested conditions, allowing for more complex specialization rules.
For example, we can have a condition that checks if a type is a vector of types that implement a specific trait, e.g.\@ \texttt{all(T = Vec<U>, U:\@ Display)}, where both \texttt{T} and \texttt{U} are generic type parameters.
Or we could have a condition that checks if a type is a reference to a type that is either \texttt{i32} or another reference to any other type, e.g.\@ \texttt{T = \&U, any(U = i32, U = \&\_)}, and so on.
All the predicates and wildcards defined earlier can be nested within each other to form more complex conditions and allow for fine-grained control over specialization.

\paragraph{Default Implementation}
In addition to specialized implementations, \mytool{} supports a default implementation that is used when no other conditions are met.
These default implementations are picked up by default by the \nameref{sec:build_script} and don't need to be explicitly marked with any macro.
Optionally, if we wanted to limit the implementations available for specialization, we could use another macro to explicitly mark the default implementations.

\paragraph{Normalization}
To ensure that the conditions are consistent and can be easily compared, \mytool{} normalizes the conditions.
This normalization process involves transforming the conditions into a canonical form, which allows for easier comparison and resolution of overlapping conditions.
The normal form chosen is a disjunctive normal form (DNF)~\cite{DNFPost21}, where each condition is represented as a conjunction of predicates.
This means that each condition is a logical OR of one or more predicatesand each predicate is a logical AND of one or more atomic conditions.
The normalization process also involves applying De Morgan's laws~\cite{DeMorgan} to push negations down to the atomic conditions, ensuring that the conditions are expressed in a consistent manner.
This representation allows for efficient checking of whether a condition is satisfied, because we can easily check each disjunctive clause separately.
For example, the condition \texttt{not(all(\mbox{T:\@ Display}, any(\mbox{T = i32}, \mbox{T = f64})))} is normalized to \texttt{any(not(\mbox{T:\@ Display}), all(not(\mbox{T = i32}), not(\mbox{T = f64})))}, so we know that that specific specialized implementation should be applied either if \texttt{T} does not implement \texttt{Display} or if \texttt{T} is neither \texttt{i32} nor \texttt{f64}.
Some particular cases could also be completely flattened during normalization, e.g.\@ \texttt{all(T = Vec<U>, any(U = i32, U = f64))} would be flattened to \texttt{any(T = Vec<i32>, T = Vec<f64>)} and hence an implementation would be generated for each of the two cases, as explained in the following paragraph.

\paragraph{Macro expansion}
Whenever possible, the conditions are evaluated at compile time and the generic parameters are replaced or constrained accordingly in all occurrences.
For example: if we have a type equality condition, the generic parameter is replaced with the specific type in all its occurrencies; if we have a trait implementation condition, the generic parameter is constrained to implement the specified trait(s).
Given that all the conditions are normalized, the macro expansion can be greatly optimized, allowing for efficient resolution of the conditions at compile time.
Each disjunct condition is expanded into its own \texttt{trait} and \texttt{impl} couple, where the new trait's name is originated by appending the conditions hash to the original trait name.
This is the reason why we previously defined this approach \textit{monomorphic specialization}.
For ease of read, in this thesis we will give each specialized trait a name that represents the underlying conditions, instead of the actual names the tool would output.
An example of this expansion is the "Specialized implementation for MyI8 or \&u8" shown in Listings~\ref{fig:conditions-definition-pre} and~\ref{fig:conditions-definition-post}, which we'll see in the following paragraph.

\paragraph{Examples}
Listings \ref{fig:conditions-definition-pre} and \ref{fig:conditions-definition-post} show an example of how to define specialized implementations using the \inlinerust{#[when(predicate)]} macro and the corresponding expansion generated by the macro.
The examples generate specializations for a struct \inlinerust{MyType} and a trait \inlinerust{Foo}, which takes a generic type parameter \inlinerust{T} and has only one method \inlinerust{foo(&self, x: T)}, as seen in Listing \ref{fig:conditions-definition-common}.

\begin{figure}[htbp!]
	\centering
	\showrust*[0.98\linewidth]{conditions-definition-common.rs}
	\captionsetup{type=listing,skip=2pt}
	\caption{Definition of the simple struct and trait used in the following listings}
	\label{fig:conditions-definition-common}
\end{figure}

\begin{figure}[htbp!]
	\centering
	\showrust*[0.98\linewidth]{conditions-definition-pre.rs}
	\captionsetup{type=listing,skip=2pt}
	\caption{Example of the conditions definition using the tool}
	\label{fig:conditions-definition-pre}
\end{figure}

\begin{figure}[htbp!]
	\centering
	\showrust*[0.98\linewidth]{conditions-definition-post.rs}
	\captionsetup{type=listing,skip=2pt}
	\caption{Example of the expansion of the conditions definition defined in Listing \ref{fig:conditions-definition-pre}}
	\label{fig:conditions-definition-post}
\end{figure}

% --------------------- Section 3 ---------------------

\section{Method Wrapping}\label{sec:method_wrapping}
To enable the selection of the appropriate specialized implementation, we need a second macro that wraps method calls.
In \mytool{}, this is accomplished using the \inlinerust{spec!} macro, which takes as first arguments the method call itself.

\paragraph{Type Annotations}
If the language supported reflection, if we had some way to inspect the type of the arguments at compile time or if this was implemented directly in the compiler, the method call alone would have been enough.
That not being the case for Rust, we need to provide additional information to the macro.
This additional pieces of information include the types of object on which the method is called and the types of the arguments passed to the method, if any.
\begin{itemize}
	\item \textbf{Type of the Object}: The type of the object on which the method is called, e.g.\@ if we have \inlinerust{x.fn(arg1, arg2)}, we need to provide the type of \texttt{x}, e.g.\@ \mbox{\texttt{MyType}}.
	\item \textbf{Types of the Arguments}: A list of the types of the arguments passed to the method, if the method takes no arguments, it can be empty or omitted, e.g.\@ if we have \inlinerust{x.fn(arg1, arg2)}, we need to provide the types of \texttt{arg1} and \texttt{arg2}, e.g.\@ \mbox{\texttt{[Arg1Type, Arg2Type]}}.
\end{itemize}
With this information, the tool can solve almost all of the specialization at compile time, generating a direct call to the appropriate specialized implementation.

\paragraph{Annotations}
The \textit{type annotations} defined above are the first two of what we call \textit{annotations}, which are additional pieces of information provided to the macro to help it resolve the specialization.
In Rust, just like we can't inspect the type of a variable at compile time, there is also no way to inspect the traits implemented by a type or to know type aliases.
To overcome these limitations, \mytool{} allows the user to provide additional annotations to the method wrapping macro (once again notice that if the language supported reflection or similar features this would not be necessary, even more so if specialization was implemented directly in the compiler).
These additional annotations include:
\begin{itemize}
	\item \textbf{Implemented Traits}: A list of traits implemented by each relevant type, e.g.\@ \mbox{\texttt{i32:\@ Display}} or \mbox{\texttt{i32:\@ Debug + Clone}}.
	\item \textbf{Type Aliases}: A list of type aliases for each relevant type, e.g.\@ \mbox{\texttt{MyInt = i32}}.
\end{itemize}
These annotations can be provided as additional arguments to the \inlinerust{spec!} macro, allowing the tool to accurately resolve the specialization even in the presence of type aliases or trait bounds in the conditions.
Considering all of this, the full syntax for the \inlinerust{spec!} macro is the following: \inlinerust{spec!{ method_call; var_type; args_types; annotations... }}, where \texttt{method\_call} is the method call, \texttt{var\_type} is the type of the object on which the method is called, \texttt{args\_types} is a list of the types of the arguments passed to the method (if any) and \texttt{annotations} are any additional annotations provided by the user.
Additionally, for languages that support class inheritance, other annotations could be added to reflect hierarchy, e.g.\@ \mbox{\texttt{T1 <:\@ T2}}.

\paragraph{Overlap}
An important aspect to consider when implementing specialization is the possibility of overlapping conditions.
When multiple specialized implementations could apply to a given method call, the tool needs a way to resolve the overlap and select the most appropriate implementation.
One way to handle this is to define a priority system based on the specificity of the conditions.
For example, a condition that checks for type equality (e.g.\@ \mbox{\texttt{T = i32}}) could be considered more specific than a condition that checks for trait implementation (e.g.\@ \mbox{\texttt{T:\@ Display}}).
In \mytool{}, this is effectively how overlaps are resolved, just like in Rust's unstable specialization feature.
When multiple conditions of the same priority level overlap, the tool raises a compile-time error, prompting the user to resolve the ambiguity by refining the conditions.
A more general and broader mechanism is the \textit{lattice} rule defined in Rust original specialization RFC~\cite{SpecializationRFC}, which instead of requiring a strict priority system, allows partial ordering given that there is also an additional implementation that covers precisely the area of overlap.
Though this approach is more flexible, it is also more complex to implement and since this is not the main priority of this thesis, we opted for the simpler priority system.
The chain of priority used in \mytool{} is as follows:
\begin{enumerate}
	\item Type Equality
	\item Trait Implementation
	\item Not Type Equality
	\item Not Trait Implementation
\end{enumerate}

\begin{figure}[htbp!]
	\centering
	\begin{flalign*}
		 & \text{T = \&i32} > \text{T = \&U} \ge \text{T = \&\_} > \text{T: Display + Debug} > \text{T: Display} > \text{not(T = \&i32)} > \\
		 & \text{not(T = \&U)} \ge \text{not(T = \&\_)} > \text{not(T: Display + Debug)} > \text{not(T: Display)}
	\end{flalign*}
	\captionsetup{type=listing,skip=2pt}
	\caption{Example of priority ordering of conditions, highest to lowest. U is a generic type parameter and hence overlaps with the wildcard unless it is constrained in some other conditions}
	\label{fig:conditions-ordering}
\end{figure}

\paragraph{Examples}
Listings \ref{fig:method-wrapping-pre} and \ref{fig:method-wrapping-post} show an example of how to wrap method calls using the \inlinerust{spec!} macro and the corresponding expansion generated by the macro.
The examples use the specialized implementations defined in Listings \ref{fig:conditions-definition-pre} and \ref{fig:conditions-definition-post}.

\begin{figure}[htbp!]
	\centering
	\showrust*[0.98\linewidth]{method-wrapping-pre.rs}
	\captionsetup{type=listing,skip=2pt}
	\caption{Example of method wrapping, using implementations defined in Listing \ref{fig:conditions-definition-pre}}
	\label{fig:method-wrapping-pre}
\end{figure}

\begin{figure}[htbp!]
	\centering
	\showrust*[0.98\linewidth]{method-wrapping-post.rs}
	\captionsetup{type=listing,skip=2pt}
	\caption{Example of the expansion of method wrapping from Listing \ref{fig:method-wrapping-pre}}
	\label{fig:method-wrapping-post}
\end{figure}

\begin{figure}[htbp!]
	\centering
	\showbashconlight*[0.98\linewidth]{method-wrapping-output.sh}
	\captionsetup{type=listing,skip=2pt}
	\caption{Example of the output printed by the program from Listing \ref{fig:method-wrapping-pre}}
	\label{fig:method-wrapping-output}
\end{figure}

% --------------------- Section 4 ---------------------

\section{Build Script}\label{sec:build_script}
The final step is to correctly handle specialization no matter the order in which the various components are defined.
Without this step, if a specialized implementation is defined after a method call site that should use it, the macro would not be able to resolve the specialization correctly.
As seen in Listing~\ref{fig:monomorphic_specialization_build_rs}, to solve this issue with our implementation of \mytool{} we simply need to add one line to the \texttt{build.rs} file.
With that, the build script analyzes the code at compile time and collects all the necessary information to handle specialization correctly.
This build script scans the code for all the \inlinerust{#[when(predicate)]} attributes and collects the specialized implementations, as well as the default implementations and the traits implemented by each relevant type.
We've also implemented a mechanism to find the \textit{conditions definition} macros even if they are imported with a different name, by analyzing the macro definitions in the code, e.g.\@ \inlinerust{use mytool::when as my_when;}.
With this information, the build script generates a metadata file that is later used by the macros to resolve the method calls correctly.

\paragraph{Alternatives}\label{subsec:alternatives}
If the language didn't support build scripts or similar mechanisms, an alternative approach could be to use some additional macros to mark the default implementations and the traits explicitly.
The first is required to identify the default implementations, while the second is needed because the trait's body is used to generate the monomorphized versions of the specialized implementations.
To handle the order of definitions without a build script or similar mechanism, the easiest way would be to have a two-pass compilation process, where the first pass collects all the necessary information and the second one generates the specialized code.
However, this approach would be less efficient and more complex to implement.
Another way would be to wrap everything in a macro that collects all the necessary information, but this would be cumbersome and would not scale well for larger codebases.

% --------------------- Section 5 ---------------------

\section{Lifetimes}\label{sec:lifetimes}
One of the main issues the Rust team faced when stabilizing specialization was handling lifetimes correctly.
Lifetimes are a core feature of Rust's ownership system and ensuring that specialized implementations respect lifetimes is crucial for maintaining safety guarantees and soundness of the type system.
It was indeed due to these challenges that specialization got stalled for a long time before being eventually marked as unstable.
When designing \mytool{}, we had to carefully consider how to handle the matter.

\paragraph{Soundness Concerns}
The main soundness concern with lifetimes and specialization arises when a specialized implementation requires, implicitly or explicitly, different lifetime constraints than those present in the default implementation.
For example, if a specialized implementation returns a reference with a shorter lifetime than the default implementation, it could lead to dangling references or other unsafe behavior.

The core problem is that at code generation time (the phase where MIR is lowered to LLVM IR and then to machine code) the compiler doesn't know about lifetimes (since it has already erased all lifetime information) and so it can't use them to choose which implementation should apply.
For example, the code shown in Listings~\ref{fig:soundness-lifetimes-explicit} and \ref{fig:soundness-lifetimes-implicit} would compile with \inlinerust{#![feature(specialization)]} but lead to unsound behavior, since typecheck (another compilation phase, which is at higher level and hence still has all lifetimes information) and code-gen would disagree about which implementation should be called by \inlinerust{x.my_method()}.
Another clear case of unsoundness of specialization with lifetimes could be if the lifetime-specialized implementation returns a different type than the default implementation, causing type-check and code-gen to disagree about the type of the returned value, leading to undefined behavior.
\begin{figure}[htbp!]
	\centering
	\showrust*[0.98\linewidth]{soundness-lifetimes-explicit.rs}
	\captionsetup{type=listing,skip=2pt}
	\caption{Example of a soundness issue with explicit lifetime constraint}
	\label{fig:soundness-lifetimes-explicit}
\end{figure}

\begin{figure}[htbp!]
	\centering
	\showrust*[0.98\linewidth]{soundness-lifetimes-implicit.rs}
	\captionsetup{type=listing,skip=2pt}
	\caption{Example of a soundness issue with implicit lifetime constraint}
	\label{fig:soundness-lifetimes-implicit}
\end{figure}

\paragraph{Workarounds}
This soundness issue was reported shortly after the specialization proposal~\cite{SpecializationSoundness} and stalled its progress for a long time.
To mitigate the issue, the Rust team proposed a workaround that allows sound specialization for a subset of use cases~\cite{MinSpecialization}.
This solution, called \inlinerust{min_specialization}, is unfortunately a breaking change from current Rust code and that's one of the reasons why it didn't get stabilized.
For example, the code in Listings~\ref{fig:soundness-lifetimes-explicit} and \ref{fig:soundness-lifetimes-implicit} would not compile with \inlinerust{#![feature(min_specialization)]}.
That is because the first (the explicit one) violates the second condition of the proposal ("it is fully generic with respect to lifetimes") by adding the \texttt{T:\@ 'static} constraint; while the second (implicit one) violates the first condition ("it relies only on implied bounds") by adding the \texttt{T:\@ MyOtherTrait} bound.

As stated earlier, the fundamental problem is that the compiler doesn't know about lifetimes at code generation time even though they could affect which implementation should be selected.
The most sensible solution to this would be in our opinion to make the compiler "smarter" by simply retaining lifetime information until code generation time (even though this would require significant changes to the compiler) or at least until monomorphization, so that lifetimes could be used to select the appropriate implementation.
This ides has already been brought up in the past but it was discarded as it was considered to have a "high cost" of implementation~\cite{SpecializationSoundness}.
The only other accusation against it is that it would be harder to reason about lifetimes if they could influence code selection, but we believe that this is not a strong argument against it, given that many modern programming languages have already accustomed us to using generic programming features without significant issues.
Hence, in \mytool{} we decided to simulate the information that the "smarter" compiler would have about lifetimes by allowing specialization based on lifetimes in \textit{conditions definition} and providing lifetime annotations in \textit{method wrapping}.

\paragraph{Conditions Definition}
To allow the user to define specialized implementations based on lifetimes, we extended the \inlinerust{#[when(predicate)]} macro to support lifetime predicates.
These predicates include:
\begin{itemize}
	\item \textbf{Type Equality}: When defining a type equality with a reference type, we can also specify lifetime constraints, e.g.\@ \mbox{\texttt{T = \&'a i32}}.
	\item \textbf{Lifetime Bound}: Check if a generic type parameter has a specific lifetime, e.g.\@ \mbox{\texttt{T:\@ 'b}}.
\end{itemize}
Other predicates could be combined with these to form more complex conditions:
\begin{itemize}
	\item \textbf{Lifetime Equality}: Check if a lifetime parameter is equal to a specific lifetime, e.g.\@ \mbox{\texttt{'a = 'static}}.
	\item \textbf{Lifetime Outlives}: Check if a lifetime parameter outlives another lifetime, e.g.\@ \mbox{\texttt{'a:\@ 'b}}.
\end{itemize}
Just like for the other conditions defined earlier, also these new conditions generate explicit constraints in the expanded implementations.

\paragraph{Method Wrapping}
To provide the necessary lifetime information to the method wrapping macro, we extended the \inlinerust{spec!} macro to accept lifetime annotations.
First of all, in \textit{type annotations} reference types can now include lifetime information, e.g.\@ \mbox{\texttt{\&'a i32}}.
Additionally, we added a new type of annotation to specify the lifetimes of each relevant type:
\begin{itemize}
	\item \textbf{Lifetime Bound}: The lifetimes of each relevant type, e.g.\@ \mbox{\texttt{T:\@ 'a}}.
\end{itemize}
Furthermore, we could add a new type of annotation to specify lifetime relationships:
\begin{itemize}
	\item \textbf{Lifetime Relationships}: Specify lifetimes equality or outlives relationships, e.g.\@ \mbox{\texttt{'a:\@ 'b}} or \mbox{\texttt{'a = 'static}}.
\end{itemize}

\begin{figure}[htbp!]
	\centering
	\showrust*[0.98\linewidth]{lifetimes-pre.rs}
	\captionsetup{type=listing,skip=2pt}
	\caption{Example of specialization using lifetimes}
	\label{fig:lifetimes-specialization-pre}
\end{figure}

\paragraph{Higher Ranked Polymorphism}
Other than first-order polymorphism (i.e.\@ generics over types), Rust also supports higher-ranked polymorphism through the use of \inlinerust{for<'a>} syntax.
This feature allows for functions to be generic over lifetimes, enabling more flexible and reusable code.
This is particularly useful in scenarios like callbacks or functional programming constructs that operate on references without being tied to a specific lifetime.
To support specialization involving higher-ranked polymorphism, we extended both macros to handle higher-ranked lifetime parameters.

An example of this can be seen in Listings~\ref{fig:higher-ranked-pre} and~\ref{fig:higher-ranked-post}, where we define a specialized implementation for a trait \inlinerust{Bar} that has a method \inlinerust{bar} accepting a higher-ranked reference.
Here we can also see how \mytool{} handles returning values from specialized implementations' functions.

\begin{figure}[htbp!]
	\centering
	\showrust*[0.98\linewidth]{higher-ranked-pre.rs}
	\captionsetup{type=listing,skip=2pt}
	\caption{Example of specialization using higher-ranked polymorphism}
	\label{fig:higher-ranked-pre}
\end{figure}

\begin{figure}[htbp!]
	\centering
	\showrust*[0.98\linewidth]{higher-ranked-post.rs}
	\captionsetup{type=listing,skip=2pt}
	\caption{Example of expansion of higher-ranked polymorphism specialization from Listing \ref{fig:higher-ranked-pre}}
	\label{fig:higher-ranked-post}
\end{figure}

% --------------------- Section 6 ---------------------

\section{Limitations}\label{sec:limitations}
While \mytool{} provides a powerful mechanism for stable specialization, it does have some limitations that should be considered:
\begin{itemize}
	\item \textbf{Boilerplate}: The need for additional annotations in method wrapping can introduce some boilerplate, especially in complex scenarios with many types and traits involved.
	\item \textbf{Compile-Time Overhead}: The use of macros and build scripts can introduce some compile-time overhead, especially in large codebases with many specialized implementations.
	\item \textbf{Lifetimes Complexity}: While the tool provides mechanisms to handle lifetimes, the added complexity may lead to confusion or misuse if not properly understood by the user.
	\item \textbf{Trait Bounds}: Since the tool generates a monomorphic version of each trait involved, the ability to involve those traits in trait bounds elsewhere in the code is lost.
	\item \textbf{Limited Predicate Support}: Although the tool supports a variety of predicates for specialization, the current set of supported predicates may not cover all language feature (e.g.\@ existential types or mutable references), potentially limiting the tool's applicability in certain cases.
	\item \textbf{Same Project Traits}: The current implementation only supports specialization for traits defined within the same crate, because we need to analyze the trait's body to generate the monomorphized versions of the specialized implementations, limiting its usability with external libraries.
	\item \textbf{Warnings and Errors}: Error messages generated by the macros can be less informative than native compiler errors, making debugging more challenging.
\end{itemize}
Despite these limitations, \mytool{} represents a significant step forward in enabling stable specialization through metaprogramming, providing developers with a powerful tool to optimize their code while adhering to Rust's safety guarantees.
More generally, the concepts presented in this chapter can be adapted to other programming languages with similar capabilities, broadening the applicability of the approach beyond Rust.
